\section{Implementation}
\label{sec:impl}

The operator ships in \sys, a governed bitemporal RDF store, as
\texttt{share\_merge} (approximately 480 lines of Rust). It exposes two entry
points. \texttt{status} loads an incoming bundle, resolves the common ancestor
through the manifest's parent pointer, and reports what each side added and
removed together with the contended slots --- without writing anything, so a
peer can be inspected before it is merged. \texttt{merge} performs the algebra
of \S\ref{sec:operator} and commits, recording both parent share identifiers on
the transaction.

Three implementation choices are worth naming because they affect what the
evaluation can claim.

\paragraph{The shapes come from the bundle, not the store.} Cardinality is read
from the shapes carried inside the share being merged, not from the local
store's copy. A peer's constraints therefore travel with its data, and merging
does not silently reinterpret their graph under our schema.

\paragraph{Conflict is reported at slot granularity, with its evidence.} A
\texttt{DecisionRecord} carries the subject, the predicate, the declared
\texttt{maxCount}, and the base, ours and theirs value sets. The person
resolving it is not told that a conflict exists and left to find it; they are
told which constraint made it one and what the three candidate states were.

\paragraph{The evaluated operator is the shipped one.} The benchmark of
\S\ref{sec:synthetic} contains its own reference implementation of every arm,
including ours, so that all eight are scored under one accounting contract. The
replay of \S\ref{sec:replay} deliberately does \emph{not}: it drives
\texttt{share\_merge::\{status,merge\}} through real bundles built by the
production share path. This matters because a benchmark's copy of an operator
can agree with a benchmark's copy of an oracle for reasons that have nothing to
do with the deployed system.
